\documentclass[10pt,a4paper]{amsart}
\usepackage[margin=19mm]{geometry}
\usepackage{amsmath,amssymb,mathtools}
\usepackage[scr=rsfs,cal=euler]{mathalfa}
\usepackage{xfrac}
\usepackage[unicode,hidelinks,bookmarks=false]{hyperref}
\newcommand{\ie}{i.e.\ }
\newcommand{\cf}{cf.\ }
\newcommand{\resp}{respectively\ }
\newcommand{\Subsection}{\subsection}
\providecommand{\nobreakdash}{\nobreak\hbox{-}}
\setlength{\emergencystretch}{2em}
\multlinegap0pt
\usepackage{bm}
\usepackage{bbm}
\usepackage{dsfont}
\usepackage{xspace}
\usepackage{cancel}
\usepackage{mathtools}
\usepackage{amsthm}
\makeatletter
\@ifundefined{theorem}{\newtheorem{theorem}{Theorem}}{}
\makeatother
\title[An Optimal Least-Square Solver For Scaled Partial-Isometric ]{An Optimal Least-Square Solver For Scaled Partial-Isometric Linear Systems}
\author{Suvendu Kar, Murugesan Venkatapathi}
\date{Source: arXiv:2507.01434v1 (CC BY 4.0)}
\begin{document}
\maketitle

\noindent\emph{Source and licence.} An Optimal Least-Square Solver For Scaled Partial-Isometric Linear Systems. Version arXiv:2507.01434v1 (CC BY 4.0), distributed under the Creative Commons Attribution 4.0 International licence, as recorded in the arXiv licence metadata for this exact version and shown on the abstract page https://arxiv.org/abs/2507.01434v1. Full rights notice: licenses/arxiv-2507.01434-CC-BY-4.0.txt.\par\smallskip

\noindent\textit{Editorial scope.} Counted unit: the Theorem whose source label is \texttt{thm:direct\_solver} in the article (source0.tex lines 248--258, source label \texttt{thm:direct\_solver}), delivered together with the article's own complete original proof of it (source0.tex lines 259--296, one \texttt{proof} environment of 38 lines). No mathematics is written, repaired or re-derived: statement and proof are the article's own text line for line. Required context retained uncounted: none. Every definition, display and label of those lines is retained unchanged. Omitted from this package and not redistributed: 1--247, 297--449. Nothing inside a retained excerpt is omitted or altered: every retained body is delivered exactly as the article's lines are written, so no omission receipt of any kind accompanies this package. Citations: the retained excerpts carry no citation command at all, so no bibliography entry of the article is transcribed and no reference is redistributed. Reference adaptation: none was needed and none was made; every reference inside the delivered unit resolves inside it, so no unresolved reference or citation is produced. Printed slips kept literal: no printed word, symbol, hypothesis or number of the counted statement or of its proof was altered. Layout and class: the article's own class is \texttt{cas-sc} with packages including natbib, cleveref, amsthm, algorithm, algpseudocode, float; the wrapper is the lane's pinned \texttt{amsart} editorial wrapper at 10pt on A4, which restates the theorem-like environments the retained text needs and adds the article's own macro definitions that the retained text needs (none), with extra wrapper packages (bm, bbm, dsfont, xspace, cancel, mathtools). The delivered numbering is the unit's own. Assets: no retained span contains a figure environment, a graphics-inclusion command, a PDF-page inclusion, a table or a TikZ picture; no image file is copied into this package, the package declares no asset dependency and no image file is redistributed with it. Licence: version arXiv:2507.01434v1 of this article is distributed under the Creative Commons Attribution 4.0 International licence, as recorded in the arXiv licence metadata for this exact version and shown on the abstract page https://arxiv.org/abs/2507.01434v1. A full rights notice is delivered at licenses/arxiv-2507.01434-CC-BY-4.0.txt. Characters: every retained excerpt is pure ASCII, so the delivered engine, XeLaTeX with its default Latin Modern text font, reports no missing character.
\begin{theorem}\label{thm:direct_solver}
For a matrix $A \in \mathbb{C}^{m\times n}$ with at least one non-zero singular value and all non-zero singular values being the same, the least square solution of the system $Ax=b$ is given by: 
\begin{equation}
x_i =
\begin{cases}
\frac{(\langle A^{(i)}, b\rangle)^*(\langle A^{(i)}, b\rangle)}{(A^{*}b)^{*}A^{*}A^{(i)}} & \text{if } b^{*}AA^{*}A^{(i)} \neq 0 \\
0 & \text{otherwise}
\end{cases}
\label{eq:a1}
\end{equation}
\end{theorem}

\begin{proof}
The case where A does not have a non-zero singular value is trivial.

Let the partial singular value decomposition of A be $A=U\Sigma V^{*}$, where $U \in \mathbb{C}^{m\times r}, \Sigma \in \mathbb{R}^{r\times r}, V \in \mathbb{C}^{n\times r}$. Let $\alpha$ be the non-zero singular values. Then the least square solution of the system of equations $Ax=b$ is given as $x= V\Sigma ^{-1}U^{*}b$ and thus 
\begin{equation}
x_i= \langle e^{(i)},x\rangle=e^{(i)^*}x=V_{(i)}\Sigma ^{-1}U^{*}b=\frac{1}{\alpha}V_{(i)}U^{*}b
\label{eq:a2}
\end{equation}
[ $e^{(i)}$ being an $n$-dimensional canonical-basis vector with the $i^{th}$ entry as 1 and all other entries as 0.]\\

Now, \begin{align}
\label{eq:a3.1}
& A^{(i)}=U\Sigma V_{(i)}^{*}\\ 
\label{eq:a3.2}
 & A^{(i)^{*}}b=V_{(i)}\Sigma U^{*}b\\
 \label{eq:a3.3}
 & A^{*}b=V\Sigma U^{*}b\\
 \label{eq:a3.4}
 & A^{*}A^{(i)}=V \Sigma ^{2}V_{(i)}^{*}
\end{align}

The condition $b^{*}AA^{*}A^{(i)} = 0$ $ \implies b^{*}U\Sigma V^{*}V\Sigma U^{*}U\Sigma V_{(i)}^{*} =0$, using the relations \eqref{eq:a3.1}- \eqref{eq:a3.4}.

\begin{align*}
 & \implies b^{*}U\Sigma ^{3} V_{(i)}^{*} =0 \\
& \implies \alpha ^3 b^*UV_{(i)}^*=0 \\
& \implies \alpha ^4 \frac{1}{\alpha}V_{(i)}U^{*}b=0\\
& \implies \alpha^4 x_i = 0, \text{ using } \eqref{eq:a2},
\end{align*}
which validates our formulation of considering $x_i=0$ when $b^{*}AA^{*}A^{(i)} = 0$ and $\alpha \neq 0$ in $\eqref{eq:a1}$.

Now, when $b^{*}AA^{*}A^{(i)} \neq 0$, the proposed solution
\begin{align*}
    x_i &= \frac{(\langle A^{(i)}, b\rangle)^*(\langle A^{(i)}, b\rangle)}{(A^{*}b)^{*}A^{*}A^{(i)}} \text{ where using relations (3)-(6) above,}\\
    &=\frac{\alpha^ 2b^*UV_{(i)}^{*}V_{(i)}U^{*}b}{\alpha^3 b^*UV_{(i)}^*}=\frac{1}{\alpha}V_{(i)}U^{*}b
\end{align*}
thus satisfying the required relation \eqref{eq:a2} for the solution.
\end{proof}

\end{document}
